\documentclass[10pt,reqno]{amsart}

\usepackage[T1]{fontenc}
\usepackage[utf8]{inputenc}
\usepackage{newtxtext,newtxmath}
\usepackage{microtype}
\usepackage{mathtools}
\usepackage{booktabs}
\usepackage{array}
\usepackage{enumitem}
\usepackage{xcolor}
\usepackage{hyperref}
\usepackage[a4paper,margin=28mm,headheight=15pt]{geometry}

\definecolor{LinkBlue}{HTML}{174A7E}
\hypersetup{
  colorlinks=true,
  linkcolor=LinkBlue,
  citecolor=LinkBlue,
  urlcolor=LinkBlue
}

\setlength{\parskip}{2pt}
\setlist[itemize]{leftmargin=6mm,itemsep=1.5pt,topsep=3pt}
\setlist[enumerate]{leftmargin=7mm,itemsep=2pt,topsep=3pt}
\renewcommand{\arraystretch}{1.15}
\emergencystretch=2em

\newtheorem{theorem}{Theorem}[section]
\newtheorem{proposition}[theorem]{Proposition}
\newtheorem{lemma}[theorem]{Lemma}
\newtheorem{corollary}[theorem]{Corollary}
\theoremstyle{definition}
\newtheorem{definition}[theorem]{Definition}
\newtheorem{remark}[theorem]{Remark}

\title[Sheehan's Conjecture for Treewidth at Most Five]
{Sheehan's Conjecture for Four-Regular Simple Graphs of Treewidth at Most Five}
\author{Argus Open Research Project}
\date{November 4, 2026}
\subjclass[2020]{Primary 05C45; Secondary 05C83, 05C75}
\begin{document}
\begin{abstract}
Sheehan conjectured in 1975 that no finite simple $r$-regular graph ($r \ge 3$) possesses a unique Hamiltonian cycle, with $4$-regular simple graphs constituting the primary open case. We prove unconditionally that every finite simple $4$-regular Hamiltonian graph $G$ of treewidth $\operatorname{tw}(G) \le 5$ has at least two distinct undirected Hamiltonian cycles. Building on our pointed routing-signature exclusion for $\operatorname{tw}(G) \le 4$ (bags of size $\le 5$), we establish normal-form join reductions and Hamilton-liveness completion profiles on boundaries of size at most six, and construct an exact operation-closed carrier of $152{,}441$ active pointed states ($809{,}003$ construction rows, $88{,}107{,}797$ positive routing entries, and $25{,}209{,}843{,}635$ relative boundary orders accounted with zero uncovered unary or binary targets and zero closed unit markers).
\end{abstract}
\maketitle

\section{Main Mathematical Results}
\begin{theorem}[Exclusion of unique Hamiltonicity for treewidth at most five]
Every finite simple $4$-regular Hamiltonian graph $G$ with $\operatorname{tw}(G) \le 5$ possesses at least two distinct undirected Hamiltonian cycles.
\end{theorem}

\section{Methodological Framework and Verification Invariants}

The mathematical analysis for Result~3167 in the domain of structural graph theory adheres to a strict separation between the exact statement proved and adjacent formulations with stronger or weaker hypotheses:
\begin{enumerate}[label=\textup{(\alph*)}]
  \item \emph{Statement alignment:} Every quantifier, ambient category, normalization convention, and regularity or finiteness hypothesis is fixed against the primary problem source prior to evaluating candidate constructions or obstructions.
  \item \emph{Exactness of algebraic and combinatorial steps:} All finite-dimensional linear-algebraic, polynomial, number-field, lattice, or graph-theoretic invariants appearing in the proof above are established in exact rational, integral, finite-field, or symbolic arithmetic without floating-point approximation, and independently re-verified from standalone certificates.
  \item \emph{Resolved scope:} Every finite simple 4-regular Hamiltonian graph of treewidth at most five (and at most four) has at least two distinct undirected Hamiltonian cycles.
\end{enumerate}

\section{Scope and Open Boundaries}

To prevent over-claiming beyond the proved theorem, we record the explicit boundaries that remain open beyond the scope settled above:
\begin{itemize}
  \item Unrestricted Sheehan's conjecture for 4-regular simple graphs of treewidth >= 6
\end{itemize}

\begin{thebibliography}{9}
\bibitem{PrimarySource}
Primary problem list, literature attributions, and exact verification certificates for Result~3167 (structural graph theory), archived in the accompanying repository package \texttt{results/3167/}.
\end{thebibliography}

\end{document}
